\subsection{Step Integrators and Change of Variable}

\begin{definition} \label{def:unit-step-function-analysis} \hlt{Unit Step Function}\\
The \hlt{unit step function} is $I:\R\to\R$ with $I(x)\defeq0$ for $x\leq0$ and $I(x)\defeq1$ for $x>0$.
\end{definition}

\begin{theorem} \label{thm:single-step-riemann-stieltjes-integral} \hlt{Single Jump Integrator}\\
Let $a<s<b$, let $f:[a,b]\to\R$ be bounded and continuous at $s$, and set $\alpha(x)\defeq I(x-s)$. Then:
\begin{enumerate}[label=(\roman*)]
\setlength{\itemsep}{0pt}
\item $f\in\mathcal{R}(\alpha)$.
\item $\int_a^b f\,d\alpha=f(s)$.
\end{enumerate}
\begin{proof}
Fix $\varepsilon>0$. Choose $\delta>0$, $(\forall x\in[a,b])\,\abs{x-s}<\delta\implies\abs{f(x)-f(s)}<\frac{\varepsilon}{2}$. Choose $x_2\in(s,\min\{b,s+\delta\})$, $\Pi\defeq\{a,s,x_2,b\}$.\\
The sole nonzero Stieltjes increment is the unit increment on $[s,x_2]$. Hence $U(\Pi,f,\alpha)-L(\Pi,f,\alpha)=\sup\limits_{[s,x_2]}f-\inf\limits_{[s,x_2]}f<\varepsilon$.
\begin{enumerate}[label=(\roman*)]
\setlength{\itemsep}{0pt}
\item The function $f$ belongs to $\mathcal{R}(\alpha)$ by Theorem \ref{thm:riemann-stieltjes-integrability-criterion}.
\item Choose tagging $t$ whose tag on $[s,x_2]$ is $s$. Then $S_\alpha(f;\Pi,t)=f(s)$; $\abs{\int_a^b f\,d\alpha-f(s)}<\varepsilon$ by Theorem \ref{thm:tagged-riemann-stieltjes-sums}(ii), $\varepsilon>0$. QED.
\end{enumerate}
\hlt{Structural Commentary}: The unit jump reduces Darboux and tagged sums to the interval containing $s$.\\
\hlt{Key Insight}: Only the jump interval carries Stieltjes mass.
\end{proof}
\end{theorem}

\begin{remark} \label{rem:stieltjes-increments-as-alpha-lengths} \hlt{Stieltjes Increments as $\alpha$-Lengths}\\
Relative to Definition \ref{def:riemann-stieltjes-integral}, the increments $\Delta\alpha_i$ are interval lengths measured by $\alpha$.\\
For $\alpha(x)=x$, they are ordinary lengths; at a jump, they concentrate weight at one point as in Theorem \ref{thm:single-step-riemann-stieltjes-integral}.
\end{remark}

\begin{theorem} \label{thm:pure-step-riemann-stieltjes-series} \hlt{Series from Pure Step Integrators}\\
Let $(c_n)_{n\geq1}\subseteq[0,\infty)$ satisfy $\sum\limits_{n=1}^{\infty}c_n<\infty$, $(s_n)_{n\geq1}$ be distinct points of $(a,b)$, $\alpha(x)\defeq\sum\limits_{n=1}^{\infty}c_nI(x-s_n)$, $f:[a,b]\to\R$ continuous.
\begin{enumerate}[label=(\roman*)]
\setlength{\itemsep}{0pt}
\item The series defining $\alpha$ converges on $[a,b]$, and $\alpha$ is increasing.
\item $f\in\mathcal{R}(\alpha)$.
\item $\int_a^b f\,d\alpha=\sum\limits_{n=1}^{\infty}c_nf(s_n)$.
\end{enumerate}
\begin{proof}
\begin{enumerate}[label=(\roman*)]
\setlength{\itemsep}{0pt}
\item If $p>q$, $(\forall x\in[a,b])\,0\leq\sum\limits_{n=q+1}^{p}c_nI(x-s_n)\leq\sum\limits_{n=q+1}^{p}c_n$. Scalar tail tends to zero, so defining series converges uniformly on $[a,b]$.\\
If $x\leq y$, then $(\forall n\geq1)\,I(x-s_n)\leq I(y-s_n)$; hence $\alpha(x)\leq\alpha(y)$.
\item The claim follows by Theorem \ref{thm:continuous-functions-riemann-stieltjes-integrable}.
\item For $N\geq1$, set $\alpha_N(x)\defeq\sum\limits_{n=1}^{N}c_nI(x-s_n)$ and $\beta_N\defeq\alpha-\alpha_N$.\\
Both are increasing by $c_n\geq0$ and monotonicity of $I$, and $f\in\mathcal{R}(\beta_N)$ by Theorem \ref{thm:continuous-functions-riemann-stieltjes-integrable}.\\
The finite identity $\int_a^b f\,d\alpha_N=\sum\limits_{n=1}^{N}c_nf(s_n)$ holds by Theorems \ref{thm:single-step-riemann-stieltjes-integral} and \ref{thm:riemann-stieltjes-integral-properties}(v).\\
Terms with $c_n=0$ vanish by the Darboux sums. Set $M\defeq\sup\limits_{x\in[a,b]}\abs{f(x)}$.\\
The identity $\abs{\int_a^b f\,d\alpha-\sum\limits_{n=1}^{N}c_nf(s_n)}=\abs{\int_a^b f\,d\beta_N}$ holds by Theorem \ref{thm:riemann-stieltjes-integral-properties}(v). Since $a<s_n<b$, $\beta_N(b)-\beta_N(a)=\sum\limits_{n>N}c_n$.\\
Hence $\abs{\int_a^b f\,d\beta_N}\leq M\sum\limits_{n>N}c_n\to0$ by Theorem \ref{thm:riemann-stieltjes-integral-properties}(iv); $\int_a^b f\,d\alpha=\sum\limits_{n=1}^{\infty}c_nf(s_n)$.
\end{enumerate}
\hlt{Structural Commentary}: Finite jump identities determine partial sums; the uniform bound controls the summable tail.\\
\hlt{Key Insight}: Each jump contributes $c_nf(s_n)$.
\end{proof}
\end{theorem}

\noindent Figure \ref{fig:riemann-stieltjes-step-integrator} depicts the first three jumps of the pure-step integrator.\\
The jump $c_i$ at $s_i$ corresponds to the summand $c_if(s_i)$ in Theorem \ref{thm:pure-step-riemann-stieltjes-series}(iii).
\begin{illustrationfigure}[H]
  \centering
  \begin{tikzpicture}[
  x=0.82cm,
  y=0.48cm,
  >=Stealth,
  every node/.style={font=\scriptsize},
  axis/.style={->, gray!70, line width=0.45pt},
  stepgraph/.style={blue!70!black, line width=0.95pt},
  jump/.style={orange!80!black, line width=0.85pt},
  point/.style={circle, fill=blue!70!black, inner sep=1.2pt},
  guide/.style={gray!45, dashed, line width=0.35pt}
]
  \draw[axis] (-0.35,0) -- (7.25,0) node[right] {$x$};
  \draw[axis] (0,-0.2) -- (0,3.0) node[above] {$\alpha(x)$};

  \coordinate (A) at (0,0.35);
  \coordinate (S1L) at (1.55,0.35);
  \coordinate (S1U) at (1.55,1.15);
  \coordinate (S2L) at (3.65,1.15);
  \coordinate (S2U) at (3.65,1.95);
  \coordinate (S3L) at (5.65,1.95);
  \coordinate (S3U) at (5.65,2.75);
  \coordinate (B) at (7.0,2.75);

  \draw[stepgraph] (A) -- (S1L);
  \draw[jump] (S1L) -- (S1U) node[midway, right] {$c_1$};
  \draw[stepgraph] (S1U) -- (S2L);
  \draw[jump] (S2L) -- (S2U) node[midway, right] {$c_2$};
  \draw[stepgraph] (S2U) -- (S3L);
  \draw[jump] (S3L) -- (S3U) node[midway, right] {$c_3$};
  \draw[stepgraph] (S3U) -- (B);

  \foreach \x/\lab in {1.55/$s_1$,3.65/$s_2$,5.65/$s_3$} {
    \draw[guide] (\x,0) -- (\x,2.85);
    \node[below] at (\x,0) {\lab};
  }

  \node[point] at (S1U) {};
  \node[point] at (S2U) {};
  \node[point] at (S3U) {};
\end{tikzpicture}

  \caption{A pure-step integrator has jump $c_i$ at $s_i$; closed lower and open upper markers record $I(0)=0$.}
  \label{fig:riemann-stieltjes-step-integrator}
\end{illustrationfigure}

\begin{theorem} \label{thm:differentiable-integrator-reduces-to-riemann-integral} \hlt{Differentiable Integrators Reduce to Riemann Integrals}\\
Let $\alpha:[a,b]\to\R$ be continuous and increasing on $[a,b]$ and differentiable on $(a,b)$.\\
Let $q\in\mathcal{R}[a,b]$ satisfy $q|_{(a,b)}=\alpha'$, and let $f:[a,b]\to\R$ be bounded. Then:
\begin{enumerate}[label=(\roman*)]
\setlength{\itemsep}{0pt}
\item $f\in\mathcal{R}(\alpha)\Longleftrightarrow fq\in\mathcal{R}[a,b]$.
\item Under these conditions, $\int_a^b f\,d\alpha=\int_a^b f(x)q(x)\,dx$.
\end{enumerate}
\begin{proof}
Set $M\defeq\sup\limits_{x\in[a,b]}\abs{f(x)}$. For partition $\Pi$, bounded $h$, set $D_\alpha(\Pi,h)\defeq U(\Pi,h,\alpha)-L(\Pi,h,\alpha)$, $D(\Pi,h)\defeq U(\Pi,h)-L(\Pi,h)$.\\
Set $E(\Pi)\defeq MD(\Pi,q)$. For tagging $t$, choose $u_i\in(x_{i-1},x_i)$ with $\Delta\alpha_i=q(u_i)\Delta x_i$ for $1\leq i\leq n$ by Corollary \ref{cor:mean-value-theorem}.\\
Then $\abs{S_\alpha(f;\Pi,t)-R(fq,\Pi,t)}\leq E(\Pi)$ by the intervalwise oscillation of $q$.\\
Hence $\abs{D_\alpha(\Pi,f)-D(\Pi,fq)}\leq2E(\Pi)$ by taking suprema over pairs of taggings.\\
The gaps $D_\alpha(\Pi,f)$ and $D(\Pi,fq)$ decrease under refinement by Theorem \ref{thm:refinement-monotonicity-riemann-stieltjes}.\\
The quantity $E(\Pi)=MD(\Pi,q)$ decreases by the same theorem with integrator $x$.\\
The relation $(\forall\delta>0)(\exists\Pi)\,E(\Pi)<\delta$ holds by Theorem \ref{thm:riemann-stieltjes-integrability-criterion} with integrator $x$.
\begin{enumerate}[label=(\roman*)]
\setlength{\itemsep}{0pt}
\item[$(i)$] ($\Rightarrow$) Fix $\varepsilon>0$. Choose $\Pi$ with $\max\{2D_\alpha(\Pi,f),4E(\Pi)\}<\varepsilon$ by the integrability criterion and refinement.\\
Then $D(\Pi,fq)\leq D_\alpha(\Pi,f)+2E(\Pi)<\varepsilon$, so $fq\in\mathcal{R}[a,b]$ by Theorem \ref{thm:riemann-stieltjes-integrability-criterion} with integrator $x$.
\item[$(i)$] ($\Leftarrow$) Fix $\varepsilon>0$. Choose $\Pi$ with $\max\{2D(\Pi,fq),4E(\Pi)\}<\varepsilon$ by the integrability criterion and refinement.\\
Then $D_\alpha(\Pi,f)\leq D(\Pi,fq)+2E(\Pi)<\varepsilon$, so $f\in\mathcal{R}(\alpha)$ by Theorem \ref{thm:riemann-stieltjes-integrability-criterion}.
\item[$(ii)$] Assume the equivalent conditions, fix $\varepsilon>0$, and set $I_\alpha\defeq\int_a^b f\,d\alpha$ and $I_q\defeq\int_a^b fq\,dx$.\\
Choose $\Pi$ with $\max\{3D_\alpha(\Pi,f),3D(\Pi,fq),3E(\Pi)\}<\varepsilon$ by the integrability criterion and refinement.\\
For any tagging $t$, each tagged sum is within $\frac{\varepsilon}{3}$ of its integral by Theorem \ref{thm:tagged-riemann-stieltjes-sums}(ii).\\
Hence $\abs{I_\alpha-I_q}<\frac{2\varepsilon}{3}+E(\Pi)<\varepsilon$ by the tagged-sum comparison. Thus $I_\alpha=I_q$.
\end{enumerate}
\hlt{Structural Commentary}: Mean-value points compare tagged sums; derivative gaps transfer integrability and identify limits.\\
\hlt{Key Insight}: The sole comparison error is the Darboux gap of $q$.
\end{proof}
\end{theorem}

\begin{theorem} \label{thm:riemann-stieltjes-change-of-variable} \hlt{Change of Variable}\\
Let $\varphi:[A,B]\to[a,b]$ be strictly increasing, continuous, and onto. Let $\alpha:[a,b]\to\R$ be increasing and $f\in\mathcal{R}(\alpha)$.\\
Set $\beta\defeq\alpha\circ\varphi$ and $g\defeq f\circ\varphi$. Then:
\begin{enumerate}[label=(\roman*)]
\setlength{\itemsep}{0pt}
\item $\beta$ is increasing on $[A,B]$.
\item $g\in\mathcal{R}(\beta)$.
\item $\int_A^B g\,d\beta=\int_a^b f\,d\alpha$.
\item Suppose $\alpha=\operatorname{id}_{[a,b]}$ and $\varphi$ is differentiable on $(A,B)$. If $q\in\mathcal{R}[A,B]$ satisfies $q|_{(A,B)}=\varphi'$, then $gq\in\mathcal{R}[A,B]$.
\item Under the hypotheses of Part (iv), $\int_a^b f(x)\,dx=\int_A^B f(\varphi(y))q(y)\,dy$.
\end{enumerate}
\begin{proof}
\begin{enumerate}[label=(\roman*)]
\setlength{\itemsep}{0pt}
\item The composition $\beta=\alpha\circ\varphi$ is increasing by the monotonicity of $\alpha$ and $\varphi$.
\item Fix $\varepsilon>0$. Choose partition $\Pi=\{x_0,\ldots,x_n\}$ with $U(\Pi,f,\alpha)-L(\Pi,f,\alpha)<\varepsilon$ by Theorem \ref{thm:riemann-stieltjes-integrability-criterion}. Set $Q\defeq\{\varphi^{-1}(x_0),\ldots,\varphi^{-1}(x_n)\}$.\\
The map $\varphi$ sends each interval of $Q$ onto the corresponding interval of $\Pi$. Hence $g$ and $f$ have the same range there and $\Delta\beta_i=\Delta\alpha_i$.\\
Hence $U(Q,g,\beta)=U(\Pi,f,\alpha)$ and $L(Q,g,\beta)=L(\Pi,f,\alpha)$. Thus $g\in\mathcal{R}(\beta)$ by Theorem \ref{thm:riemann-stieltjes-integrability-criterion}.
\item Upper-sum, lower-sum families coincide by bijection $\Pi\mapsto\varphi^{-1}(\Pi)$ from Part (ii). Hence $\int_A^B g\,d\beta=\int_a^b f\,d\alpha$ by Part (ii),  $f\in\mathcal{R}(\alpha)$.
\item Here $\beta=\varphi$ and $g\in\mathcal{R}(\varphi)$ by Part (ii). Hence $gq\in\mathcal{R}[A,B]$ by Theorem \ref{thm:differentiable-integrator-reduces-to-riemann-integral}(i).
\item The equality $\int_a^b f(x)\,dx=\int_A^B g\,d\varphi=\int_A^B f(\varphi(y))q(y)\,dy$ holds by Part (iii) and Theorem \ref{thm:differentiable-integrator-reduces-to-riemann-integral}(ii).
\end{enumerate}
\hlt{Structural Commentary}: Pullback identifies Darboux sums; derivative reduction replaces $d\varphi$ by $q\,dy$.\\
\hlt{Key Insight}: Corresponding intervals have identical function ranges and integrator increments.
\end{proof}
\end{theorem}
